% Tipografía: TeX Gyre Pagella (clon libre de Palatino, licencia GUST).
%
% Las fuentes se piden por NOMBRE DE ARCHIVO y no por nombre de familia: los
% .otf viven en el árbol texmf de TeX Live y XeTeX los encuentra por kpathsea,
% mientras que buscarlos por nombre exige que fontconfig los tenga registrados
% —cosa que no ocurre ni en una TinyTeX recién instalada ni en el runner de CI—.
% Así el PDF compila igual en cualquier máquina con TeX Live, sin instalar
% fuentes en el sistema.
\setmainfont{texgyrepagella}[
  Extension      = .otf,
  UprightFont    = *-regular,
  BoldFont       = *-bold,
  ItalicFont     = *-italic,
  BoldItalicFont = *-bolditalic,
]
\setsansfont{texgyreheros}[
  Extension      = .otf,
  UprightFont    = *-regular,
  BoldFont       = *-bold,
  ItalicFont     = *-italic,
  BoldItalicFont = *-bolditalic,
  Scale          = MatchLowercase,
]
\setmonofont{texgyrecursor}[
  Extension      = .otf,
  UprightFont    = *-regular,
  BoldFont       = *-bold,
  ItalicFont     = *-italic,
  BoldItalicFont = *-bolditalic,
  Scale          = MatchLowercase,
]

% Matemática en la misma familia que el texto.
\usepackage{unicode-math}
\setmathfont{texgyrepagella-math.otf}

\usepackage{microtype}   % con xelatex: protrusión
